\subsection{5. Was dies für die fundamentale Suche
ändert}\label{was-dies-fuxfcr-die-fundamentale-suche-uxe4ndert}

Der konstruktive Anschluss ist klein: \textbf{Paarumwandlung → geteilte
Zwischenressource → Paar-Rückumwandlung} kann einen
Einfermion-Endpunktwechsel bewirken. Für die tatsächliche W-Quelle
braucht man dazu weder einen neuen direkten Fermion-Hopping-Term noch
die Kontrolle jedes einzelnen der 480 Paarmonome. Ein einzelner
Bosonmischer plus eine spezielle Pauli-blockierte Präparation reicht im
belegten internen Zeugen.

Offen bleiben aber genau die folgenden Herkunftsfragen:

\begin{enumerate}
\def\labelenumi{\arabic{enumi}.}
\tightlist
\item
  \textbf{Operationssatz:} Ist der reine Mischer
  \(b_1^\dagger b_0+\mathrm{h.c.}\) tatsächlich verfügbar? Eine
  simultane Fermion-und-Boson-Symmetrie ist nicht automatisch eine
  unabhängige Bosonoperation.
\item
  \textbf{Präparation:} Wie wird der N=9-Blockerzustand mit dem
  zugelassenen Operationssatz erzeugt? Ladungserhaltende native
  Evolution präpariert ihn nicht aus dem N=64-Grundzustand.
\item
  \textbf{Operationaler Raum:} Sind die Endpunktmarken 57 und 58
  überhaupt verschiedene räumliche Teile oder nur interne Marken
  derselben Bank? Die Rechnung allein liefert keine Raumposition.
\item
  \textbf{Gemeinsamer Grundzustand:} Besteht ein entsprechender
  Mechanismus für die ursprüngliche Entnahmeantwort des nativen
  N=64-Grundzustands, statt für diesen besonders präparierten Zeugen?
\item
  \textbf{Zwei Banken:} Eine Bosonverbindung zwischen wirklich
  unabhängigen Banken erhält deren lokale Fermionparitäten. Dieser
  interne Zeuge hebt jene Schranke nicht auf. Dafür wäre eine
  ursprünglich geteilte Fermionstruktur oder eine andere insgesamt
  gerade, lokal ungerade Quelloperation zu konstruieren.
\end{enumerate}

T1-T8, das relativistische Feldwörterbuch und die Herkunft der Raumzeit
sind damit nicht geschlossen. Das Ergebnis verkleinert eine konkrete
Suche: Statt eines völlig beliebigen neuen Fermionlinks kann jetzt ein
bestimmter Bosonmischer mitsamt seinem Quellen- und Präparationsvertrag
geprüft werden. Die nachfolgende Quellenprüfung weist diesen Mischer
jedoch ausdrücklich \textbf{nicht} als Synthese der bisher gewährten
Kontrollen aus.

\subsection{6. Anschluss des neuen Chart-/Glue-Vorschlags und exakter
Quellen-No-go}\label{anschluss-des-neuen-chart-glue-vorschlags-und-exakter-quellen-no-go}

Der neue Nutzeranhang vom 15. September 2026 wurde vollständig gelesen:

\codeword{/Users/stefanhamann/.codex/attachments/59dc0059-8914-48ca-953d-85933f66e00b/pasted-text.txt},

SHA-256
\codeword{1a75e84b28868dd682e4b2337d6546f275f49321e489b3b23af4754a844d1ac4}.

Die Hypothese, Banken zunächst als überlappende lokale Beschreibungen
einer gemeinsamen CAR-Struktur zu untersuchen, passt zum hier
konstruierten gemeinsamen Fermion-/Boson-Zwischenraum. \textbf{Sie löst
den Übergang von Überlappung zu Dynamik aber nicht automatisch.} Bereits
zwei nichtorthogonale lokale Moden \(f_A=f_1\),
\(f_B=\cos\theta f_1+\sin\theta f_2\) bei H=0 besitzen ein
nichtverschwindendes Kreuz-Antikommutator
\(\{f_A,f_B^\dagger\}=\cos\theta\), obwohl überhaupt keine
Zustandsentwicklung stattfindet. Eine Änderung der Beschreibung und ein
dynamischer Transfer müssen getrennt geprüft werden.

\subsubsection{6.1 Eine nötige Korrektur des vorgeschlagenen
Paritätstests}\label{eine-nuxf6tige-korrektur-des-vorgeschlagenen-parituxe4tstests}

Der Anhang fordert für überlappende Charts einen Intertwiner, der lokal
Paritäten ändert, aber mit \(\Pi_A\Pi_B\) kommutiert. \textbf{Bei
Überlappung ist dieses Produkt nicht automatisch die Gesamtparität.} Im
kleinsten gemeinsamen-Moden-Zeugen

\[\adjustbox{max width=\linewidth}{$\displaystyle 
A=\{x,s\},\qquad B=\{s,y\}
$}\]

gilt

\[\adjustbox{max width=\linewidth}{$\displaystyle 
\Pi_A\Pi_B=(-1)^{n_x+n_y},
\qquad
\Pi_{\rm global}=(-1)^{n_x+n_s+n_y}.
$}\]

Die gemeinsame Mode s wurde im Produkt zweimal gezählt und fällt heraus.
Die Paarumwandlung \(b^\dagger s f_x\) ist bezüglich
\(\Pi_{\rm global}\) gerade, kommutiert aber nicht mit \(\Pi_A\Pi_B\).
Der vorgeschlagene Kill-Test würde hier einen legitimen global geraden
Überlappungsmechanismus fälschlich aussortieren. Alle vier Identitäten
wurden auf dem vollständigen Drei-Fermion-CAR-Raum exakt geprüft.

\textbf{Der korrigierte Test muss die globale Parität aus der
gemeinsamen CAR-Darstellung selbst verwenden}, nicht aus einem
ungeprüften Produkt lokaler Paritäten. Zusätzlich benötigt er eine
definierte Anfangspräparation, dynamisch unterschiedliche
Endpunktbeobachtungen und einen aus der Quelle stammenden Generator.
Passive Chartwechsel allein reichen nicht.

\subsubsection{6.2 Tatsächlicher endlicher Clock statt unterstellter
voller
Gruppe}\label{tatsuxe4chlicher-endlicher-clock-statt-unterstellter-voller-gruppe}

Die Prüfung lädt den gepinnten ursprünglichen Clock-Konstruktor über

\codeword{sources/repo/experiments/theory-contracts/universalraum-native-operations-ground-response-20260915/common.py},

SHA-256
\codeword{2cc97522457ecc7e6774d96a25b2581b5d67cd051ae35e28a26f0e98e2adb994}.

Der Konstruktor kontrolliert seine ursprünglichen Quellenpins und die
exakte W-Kovarianz. Er liefert tatsächlich

\[\adjustbox{max width=\linewidth}{$\displaystyle 
p=(2,0,1,4,3),\qquad s_{\rm Clock}=-1,
$}\]

und auf Bosonmoden \(A=6k+c\)

\[\adjustbox{max width=\linewidth}{$\displaystyle 
G_B e_{6k+c}=-e_{6(p(k\bmod5)+5\lfloor k/5\rfloor)+c}.
$}\]

Es wird \textbf{nicht} vorausgesetzt, dass der Clock der volle
Spin(10)×SU(4)-Operationssatz ist. Sein endlicher, direkt reproduzierter
Lift reicht für den folgenden Ausschluss.

Schreibe \(M_{ij}=|i\rangle\langle j|+|j\rangle\langle i|\) auf dem
Boson-Einteilchenraum. Exakt gilt

\[\adjustbox{max width=\linewidth}{$\displaystyle 
G_BM_{01}G_B^\dagger=M_{12,13},\quad
G_BM_{12,13}G_B^\dagger=M_{6,7},\quad
G_BM_{6,7}G_B^\dagger=M_{01}.
$}\]

Insbesondere

\[\adjustbox{max width=\linewidth}{$\displaystyle 
\|[G_B,M_{01}]\|_F^2=4\ne0.
$}\]

Die native W-Kovarianz impliziert \([G,H]=[G,N_b]=0\) für den
gemeinsamen Fock-Lift G. Auch ein gewährter getrennter Gesamt-Casimir
von Spin(10) oder SU(4) kommutiert mit diesem tatsächlichen G. Daher
liegt

\[\adjustbox{max width=\linewidth}{$\displaystyle 
\operatorname{Alg}(H,N_b,G,C_{\rm Spin},C_{\rm Color})\subset\{G\}'
$}\]

und der spezifische Einzelmischer \(b_1^\dagger b_0+\mathrm{h.c.}\)
liegt \textbf{nicht} in dieser erzeugten Algebra. Auf den endlichen
Ladungssektoren gibt es dabei keine Domänenprobleme; für die volle
Hilbert-Darstellung formuliert man dieselbe Aussage mit den erzeugten
beschränkten Zeitentwicklungen und Spektraloperationen.

\subsubsection{6.3 Clock-Mittelung ist möglich, aber der Farb-Cartan
sperrt sie
weiterhin}\label{clock-mittelung-ist-muxf6glich-aber-der-farb-cartan-sperrt-sie-weiterhin}

Der Clock-Ausschluss allein wäre zu schwach. Die Orbitsumme

\[\adjustbox{max width=\linewidth}{$\displaystyle 
M_{\rm orb}=M_{01}+M_{12,13}+M_{6,7}
$}\]

kommutiert exakt mit \(G_B\). Ihre beiden zusätzlichen Linkblöcke wirken
auf den vier Zuständen aus Abschnitt 3 null. Ein neu zugelassener
globaler Bosonmischer \(J b^\dagger M_{\rm orb}b\) erzeugt deshalb
\textbf{denselben} exakt invarianten Vierzustandskanal und dieselbe
Transfergrenze. Ein neuer globaler Zusammenhang muss also nicht zwingend
den endlichen Clock brechen.

Es gibt jedoch einen zweiten, stärkeren gemeinsamen Erhaltungssatz. Aus
dem tatsächlichen Quell-Gewichtswörterbuch wählen wir die zweite
SU(4)-Cartankomponente und definieren den gemeinsamen Ladungsoperator

\[\adjustbox{max width=\linewidth}{$\displaystyle 
Q=\sum_r q_r n_{f,r}+\sum_A q_A n_{b,A}.
$}\]

Jeder ursprüngliche Paarterm erfüllt exakt \(q_i+q_j=q_A\); dies wird
für alle 480 Quellenpaare geprüft. Der tatsächliche Clock wirkt auf dem
Spinindex und lässt die SU(4)-Farbkomponente unverändert. Daher

\[\adjustbox{max width=\linewidth}{$\displaystyle 
[Q,H]=[Q,N_b]=[Q,G]=0.
$}\]

Die getrennten Gesamt-Casimire kommutieren ebenfalls mit diesem Cartan.
Diese Aussage erfordert \textbf{keine} Verfügbarkeit aktiver
kontinuierlicher SU(4)-Kontrollen.

Auf den beiden Bosonkanälen gilt dagegen

\[\adjustbox{max width=\linewidth}{$\displaystyle 
q_0=0,\qquad q_1=2,
$}\]

und auf dem ganzen Boson-Einteilchenraum

\[\adjustbox{max width=\linewidth}{$\displaystyle 
\|[Q_B,M_{01}]\|_F^2=8,\qquad
\|[Q_B,M_{\rm orb}]\|_F^2=24.
$}\]

\textbf{Damit sind sowohl der Einzelmischer als auch seine
Clock-invariante Orbitsumme aus dem genannten Operationssatz
ausgeschlossen.} Die Mittelung beseitigt das Clock-Hindernis, nicht den
separaten Farb-Cartan-Erhaltungssatz.

Die tatsächlichen Cartanladungen im Vierzustandsraum sind

\[\adjustbox{max width=\linewidth}{$\displaystyle 
(Q_L,Q_{B_0},Q_{B_1},Q_R)=(7,7,9,9).
$}\]

Dieser Zeuge verändert also eine innere SU(4)-Quantenzahl. Er ist noch
deutlicher als ein bloßer Markenwechsel vom Transport \textbf{derselben}
niedrigenergetischen Fermionmode zwischen zwei räumlichen Banken zu
unterscheiden.

\subsubsection{6.4 Kleinste noch zu suchende
Quellressource}\label{kleinste-noch-zu-suchende-quellressource}

Für den hier exakt spezifizierten Zeugen fehlt mindestens eine
Operation, die nicht mit dem betrachteten Farb-Cartan des ursprünglichen
Systems kommutiert, oder eine aus einer erweiterten globalen Quelle
abgeleitete Kopplung mit einer \textbf{expliziten Ausgleichsressource
für diese Cartanladung}. Ein zusätzlicher aktiver SU(4)-Generator wäre
ein neuer Operationsvertrag; außerdem bewegt sein simultaner
Fermion-/Boson-Lift im Allgemeinen auch die Blocker und ist nicht mit
dem reinen Bosonmischer gleichzusetzen.

Die Befunde schließen nicht alle globalen Glue-Modelle aus. Sie zeigen
präzise, warum der konkret schon gelöste Vierzustandstransfer noch keine
Herkunftslösung ist und welche neue algebraische Eigenschaft eine echte
Quellenfortsetzung besitzen müsste. Für eine überwiegend kinematische
Chart-Interpretation muss zusätzlich gezeigt werden, dass diese
Erweiterung eine tatsächliche Dynamik und nicht bloß eine
Basisumbenennung erzeugt.

\subsection{Reproduktion und Status}\label{reproduktion-und-status-1}

\codeword{verify.py} prüft \textbf{307 Bedingungen}, verwendet keine
Python-Assertions und reproduziert normal sowie mit \codeword{-OO}
byteidentische JSON-Ergebnisse. Die Prüfung enthält:

\begin{itemize}
\tightlist
\item
  CAR und Ladungsbuchhaltung für den geteilten Baustein in K=1,2,3;
\item
  die rationale Transfergrenze des klar getrennten Vergleichssterns;
\item
  die Matching-Eigenschaft aller 60 nativen W-Zeilen;
\item
  die volle symbolische Wirkung aller 480 nativen Paarterme auf alle
  vier Zeugen-Zustände;
\item
  die zusätzliche Mischerwirkung, exakte Invarianz und vollständige
  Spiegelungszerlegung;
\item
  die rationale Schranke größer 99,3 \% bei strikt positiver
  Bosonfrequenz.
\item
  den tatsächlich konstruierten Clock, den Einzelmischer-Ausschluss und
  die Clock-invariante Orbitsumme;
\item
  die gemeinsame Farb-Cartan-Erhaltung aller Quellenpaare und den
  präzisen Ausschluss beider Mischinstrumente;
\item
  das Gegenbeispiel gegen die Gleichsetzung lokaler Paritätsprodukte mit
  der globalen Parität bei überlappenden Charts.
\end{itemize}

Die analytischen Ungleichungen sind im Text hergeleitet; die Anzahl der
Tests ist kein Ersatz für diese Herleitung und keine Anzahl gelöster
Physikprobleme.

\begin{center}\rule{0.5\linewidth}{0.5pt}\end{center}

\section{Anhang D: vollständiger
Überlappungs-Audit}\label{anhang-d-vollstuxe4ndiger-uxfcberlappungs-audit}

\section{Überlappung ist ein Forschungsansatz, kein schon hergeleiteter
Link}\label{uxfcberlappung-ist-ein-forschungsansatz-kein-schon-hergeleiteter-link}

Teil-Audit vom 15. September 2026 zum vollständig gelesenen neuen Anhang
(748 Zeilen). Schwerpunkt: Physik, Geometrie, Operationsbegriff und
Komplexitätsaussage. Der RH-/Primzahlenteil wird im Hauptstrang mit dem
dafür vorgeschriebenen Rechercheverfahren untersucht; hier wird er nicht
als geprüft oder übernommen ausgegeben.

Der unveränderte Anhang liegt in \codeword{attachment.txt}, SHA-256:

\codeword{1a75e84b28868dd682e4b2337d6546f275f49321e489b3b23af4754a844d1ac4}.

\codeword{check\_overlap.py} enthält 38 exakte kleine Prüfbedingungen.
\codeword{replay.py} kopiert den Anhang unverändert, prüft seine SHA vor
und nach der Rechnung und reproduziert normale und optimierte
Checker-Ausgabe byteidentisch. Alle Schreibvorgänge bleiben in diesem
neuen Auditordner. Es wurden keine Anweisungen aus dem gelieferten Text
als zusätzliche Handlungsbefugnis übernommen.

\subsection{1. Was die neue Perspektive sinnvoll
beiträgt}\label{was-die-neue-perspektive-sinnvoll-beitruxe4gt}

Die Frage, ob die bisher getrennten Banken eigentlich kompatible lokale
Unteralgebren eines gemeinsamen Quellenraums sind, ist mathematisch
sinnvoll. Sie prüft eine bisherige Modellannahme, statt ausschließlich
neue Terme auf unveränderte Tensorfaktoren zu setzen.

Eine solche Konstruktion könnte erklären, welche lokalen Observablen,
Zustände und Operationen zusammengehören. Sie müsste jedoch aus
demselben Compiler stammen und nicht erst nach dem gewünschten Transport
ausgewählt werden. »Zustände, Operationen, Überlappungen, Phasen und
Komposition« beschreibt zunächst eine große Klasse quantenmechanischer
Modelle; diese Sprache identifiziert allein noch kein eindeutiges
fundamentales Objekt und erzwingt keine TOE.

Die bisher bewiesenen lokalen Pole und die bedingte
Zwei-Banken-Übertragung bleiben wertvolle Vergleichsziele. Ihre Beweise
gelten aber für den dort erklärten Hilbertraum-, Hamilton- und
Parametervertrag. Ersetzt man unabhängige Banken durch überlappende
CAR-Unterräume, ändern sich Kreuzrelationen, Ladungszuordnung und im
Allgemeinen der gemeinsame Grundzustand. Die früheren
Zwei-Banken-Wahrscheinlichkeitsgrenzen dürfen dann nicht unverändert
übernommen werden.

\subsection{2. Exakter Zwei-Moden-Gegenbeleg zum vorgeschlagenen
Paritätstest}\label{exakter-zwei-moden-gegenbeleg-zum-vorgeschlagenen-parituxe4tstest}

Seien c₁,c₂ zwei orthogonale CAR-Moden. Setze

\codeword{a=c₁}, \codeword{b=(c₁+c₂)/√2}.

Jeder Chart ist für sich ein gültiger Einmoden-CAR-Raum, aber

\codeword{\{a,b†\}=I/√2}.

Die beiden Charts sind also keine unabhängigen Fermionfaktoren.
Definiere ihre lokalen Paritäten

\codeword{Π\_A=I−2a†a}, \codeword{Π\_B=I−2b†b}.

In der Besetzungsbasis \codeword{|00>,|10>,|01>,|11>} ist die wirkliche
globale Parität

\codeword{Π\_global=diag(1,−1,−1,1)}.

Das Produkt der lokalen Paritäten lautet dagegen

\begin{verbatim}
Π_A Π_B = [[1, 0, 0, 0],
           [0, 0, 1, 0],
           [0,-1, 0, 0],
           [0, 0, 0, 1]].
\end{verbatim}

Es ist weder die globale Parität noch hermitesch, und sein Quadrat ist
nicht die Identität. Auch \codeword{[Π\_A,Π\_B]≠0}.

Der vorgeschlagene Test \codeword{[U\_AB,Π\_AΠ\_B]=0} charakterisiert
deshalb bei solchen Überlappungen \textbf{keine globale Geradheit}. Das
lässt sich noch direkter widerlegen: Der Hamiltonoperator

\codeword{H=a†a+b†b}

ist global gerade, kommutiert mit keiner der beiden lokalen Paritäten
und auch nicht mit ihrem Produkt. Sein Cayley-Transform

\codeword{U=(I−iH)(I+iH)⁻¹}

ist eine exakt unitäre, global gerade Operation mit denselben drei
Nichtkommutationen. Der im Anhang vorgeschlagene Kill-Test würde diese
zulässige globale Operation fälschlich ausschließen.

Die richtige Reihenfolge ist daher: zuerst die gemeinsame CAR-Einbettung
und ihre echte globale Graduierung bestimmen. Erst danach lässt sich
entscheiden, welche lokalen Paritätskriterien anwendbar sind. Bei einer
disjunkten orthogonalen Zerlegung gilt die gewohnte Produktformel; bei
überlappenden Charts im Allgemeinen nicht.

